\chapter{Координатная алгебра M300 и корректность комплекса одноформ}

Предыдущая глава получила положительный гессиан на заранее выделенном
цепном секторе. Чтобы назвать его полным физическим гессианом, требуется
вычислить представленные одноформы. Однако для этого одного операторного
носителя и одного следа недостаточно.

Необходимо различать
\begin{equation}
 \mathcal B_{\mathrm{par}}=B(\mathcal H_{\mathrm{par}})
 \simeq M_{300}(\mathbb C)
\end{equation}
и координатную алгебру спектральной геометрии
\begin{equation}
 \mathcal A_{\mathrm{coord}}
 \xrightarrow{\ \pi\ }
 B(\mathcal H_{\mathrm{par}}).
\end{equation}
Первая фиксирует нормированный матричный след. Вторая вместе с оператором
Дирака определяет пространство
\begin{equation}
 \Omega_D^1(\mathcal A_{\mathrm{coord}})
 =\left\{
 \sum_j\pi(a_j)[D,\pi(b_j)]
 \right\}.
 \label{eq:v5-m300-oneforms-definition}
\end{equation}
В существующей склейке полная алгебра $\mathcal A_{\mathrm{coord}}$, её
левое и противоположное представления на всех 300 состояниях и блоки $D$
между опорой и цепью ещё не заданы.

\section{Почему полную матричную алгебру нельзя объявить координатной}

Рассмотрим попытку положить
\begin{equation}
 \mathcal A_{\mathrm{coord}}=M_n(\mathbb C),
 \qquad \mathcal H=\mathbb C^n
\end{equation}
в определяющем представлении. Коммутант этого представления равен
\begin{equation}
 \pi(M_n(\mathbb C))'=\mathbb C I_n.
\end{equation}
Условие нулевого порядка вещественной спектральной тройки требует
\begin{equation}
 [\pi(a),J\pi(b)J^{-1}]=0
 \qquad \text{для всех }a,b.
\end{equation}
Следовательно, образ противоположной алгебры обязан лежать в скалярах и
не может быть верным для $M_n(\mathbb C)$ при $n>1$.

Эквивалентно, верный левый--правый бимодуль над
$M_n(\mathbb C)$ является модулем над
\begin{equation}
 M_n(\mathbb C)\otimes M_n(\mathbb C)^{\mathrm{op}}
 \simeq M_{n^2}(\mathbb C)
\end{equation}
и имеет комплексную размерность не менее $n^2$. При $n=300$ это
$90\,000$, а не 300.

\begin{theorem}[Запрет определяющего представления M300]
Носитель $\mathbb C^{300}$ не допускает верного левого и правого действий
$M_{300}(\mathbb C)$, удовлетворяющих условию нулевого порядка. Поэтому
$M_{300}(\mathbb C)$ в текущей конструкции может быть алгеброй всех
операторов и носителем следа, но не координатной алгеброй вещественной
спектральной тройки в определяющем представлении.
\end{theorem}

Это стандартное следствие бимодульного описания конечных спектральных
троек и диаграмм Краевского; оно не является новым физическим запретом.

\section{Взрыв пространства одноформ без вещественной структуры}

Даже если временно отбросить противоположное действие, полная матричная
алгебра не сохраняет малый цепной сектор.

\begin{proposition}[Одноформы полной матричной алгебры]
Пусть $D$ --- нескалярный самосопряжённый оператор на $\mathbb C^n$.
Тогда в определяющем представлении
\begin{equation}
 \Omega_D^1(M_n(\mathbb C))=M_n(\mathbb C).
\end{equation}
\end{proposition}

Действительно, в собственном базисе $D$ для любого столбца $q$ можно
выбрать $j$ с $\lambda_j\ne\lambda_q$. Тогда
\begin{equation}
 E_{pj}[D,E_{jq}]
 =(\lambda_j-\lambda_q)E_{pq},
\end{equation}
поэтому получаются все матричные единицы $E_{pq}$.

Для $n=300$ это означает:
\begin{equation}
 \dim_{\mathbb C}\Omega_D^1=90\,000,
 \qquad
 \dim_{\mathbb R}(\Omega_D^1)_{\mathrm{sa}}=90\,000.
\end{equation}
Проективная калибровочная алгебра имела бы размерность $89\,999$.
Это не одиннадцать переменных $(X,\Phi)$ и не группа
$SO(3)\times SU(5)$.

\section{Точный свидетель неединственности исчисления}

Рассмотрим четырёхвершинный сокращённый носитель: одна изолированная
опорная вершина и цепь $L-G-R$. Оператор Дирака и нормированный матричный
след оставим одними и теми же.

Если координатная алгебра равна $\mathbb C^4$ и действует диагонально, то
представленные одноформы имеют комплексный ранг четыре: это две стрелки
цепи и их сопряжения. Ранг смешивания опоры с цепью равен нулю.

Если на том же носителе взять $M_4(\mathbb C)$, то ранг одноформ равен
шестнадцати, а ранг опорно-цепного смешивания --- шести. Таким образом,
\begin{equation}
 (\mathcal H,D,\tau)
 \quad\text{не определяют}\quad
 \Omega_D^1
\end{equation}
без выбора $\mathcal A_{\mathrm{coord}}$ и её представления.

Машинный аудит одновременно проверяет:
\begin{equation}
 4\ne16,
 \qquad 0\ne6,
 \qquad \dim M_4(\mathbb C)'=1.
\end{equation}

\section{Следствие для действия Ходжа}

Положительный результат предыдущей главы сохраняется как точное
утверждение на ограниченном цепном исчислении:
\begin{equation}
 S_H=\tau_{300}([d,d^\dagger]-h)^2
\end{equation}
имеет устойчивый минимум по $(X,\Phi)$. Но теперь нельзя утверждать, что
этот гессиан является ограничением уже определённого полного гессиана
$M_{300}$. Полного исчисления пока нет.

Это не опровержение самого действия. Это исправление уровня утверждения:
единственный след получен, а координатная геометрия и пространство полей
ещё не получены.

\begin{nogocriterion}[Корректность полного гессиана]
Полный гессиан и фактор БВ/БРСТ нельзя вычислять, пока заранее и независимо
от желаемого спектра не заданы
\begin{equation}
 (\mathcal A_{\mathrm{coord}},\pi,\mathcal H_{300},D,J,\gamma)
\end{equation}
на опорном, цепном и общем 45-мерном секторах. Объявление
$\mathcal A_{\mathrm{coord}}=M_{300}(\mathbb C)$ в определяющем
представлении считается провалом по условию нулевого порядка и по числу
полей.
\end{nogocriterion}

\section{Вердикт}

\begin{protocol}
\begin{enumerate}
 \item $M_{300}$ как единый носитель следа:
 \textbf{сохранён};
 \item $M_{300}$ как координатная алгебра на $\mathbb C^{300}$:
 \textbf{исключён};
 \item полный комплекс одноформ из существующих данных:
 \textbf{не определён};
 \item опорно-цепные смешивания:
 \textbf{зависят от ещё не выбранной алгебры};
 \item положительный связующий гессиан:
 \textbf{сохранён условно на цепном исчислении};
 \item полный фактор БВ/БРСТ:
 \textbf{пока не имеет однозначной постановки};
 \item физическое замыкание:
 \textbf{не достигнуто}.
\end{enumerate}
\end{protocol}

Следующий гейт должен построить минимальную координатную алгебру, её
верный бимодуль, $J$, градуировку и полный носитель оператора Дирака на
$\mathbb C^{300}$. Только после этого разрешено вычислять
$\Omega_D^1$, идеал мусора второй степени и фактор БВ/БРСТ.

Машинная проверка и сертификат сохранены в
\path{s2t_v5_m300_coordinate_algebra_wellposedness_gate.py} и
\path{s2t_v5_m300_coordinate_algebra_wellposedness_gate_results.json}.

Литературная граница сверена с классификацией конечных спектральных троек
Краевского (arXiv:hep-th/9701081), бимодульным описанием Пашке и Ситаржа
(arXiv:q-alg/9612029) и построением БВ-комплекса конечных спектральных
троек Изеппи (arXiv:2410.11823).